% Extracted from the compiled manuscript.
% Source label: tab:micro_subspace_pooled

\begin{table}[!b]
\centering
\caption{Supplementary continuous micro-subspace model ladder.}
\label{tab:micro_subspace_pooled}
\scriptsize
\setlength{\tabcolsep}{2.4pt}
\begin{tabular}{lrrrrrrr}
\toprule
Duration & \shortstack{M0\\\(R^2\)} & \shortstack{M1\\\(R^2\)} &
\shortstack{M2\\\(R^2\)} & \shortstack{Macro\\\(\Delta R^2\)} &
\shortstack{Micro\\\(\Delta R^2\)} &
\shortstack{Micro 95\%\\interval} & Significant\\
\midrule
30 s & \(-0.0886\) & 0.0123 & 0.0340 & 0.1010 & 0.0216 & [0.0118, 0.0299] & Yes\\
45 s & \(-0.0754\) & 0.0380 & 0.0648 & 0.1134 & 0.0268 & [0.0175, 0.0354] & Yes\\
60 s & \(-0.0672\) & 0.0541 & 0.0855 & 0.1213 & 0.0315 & [0.0211, 0.0412] & Yes\\
\bottomrule
\end{tabular}
\begin{minipage}{0.94\linewidth}\footnotesize
M0 is persistence, M1 adds the macro dial, and M2 adds the first continuous
principal component of the within-day residual tangent space.  Results are
pooled out-of-sample estimates under LE geometry.  The micro increment is
supplementary and post hoc.
\end{minipage}
\end{table}
